% Source: external artifact root / figure 2 / intro_comparison.tex
% Source SHA-256: 76467a9a8901131c0b1ccc885ff8ea68e6063a6bff68b33a2f4ffd85035e266c
% Only image paths below were changed to images/figure_2a.pdf, 2b.pdf, 2c.pdf.
% Put generated panel PDFs in the manuscript's images/ directory.
% The supplied captions, bibliography keys and theorem/macro dependencies are
% preserved verbatim; this snippet does not provide their definitions.
% This supplied caption is a different draft from the submitted PDF caption.
% Only the three graphic assets are verified against the submitted Figure 2.
% Preamble: \usepackage{graphicx,subcaption,natbib}
% The manuscript is assumed to define \mcA and \unerring.
% Replace TODO-global-analyses with your bibliography keys.
% Exact panel versions: allspace v3; ellipsoid v7; intersection v8 B.
% All three selected panels show -\nabla L, with no flow paths.

\begin{figure}[]
    \centering
    \begin{subfigure}[t]{0.31\linewidth}
        \centering
        \vspace{0pt}
        \includegraphics[width=\linewidth]{images/figure_2a.pdf}
        \caption{Conditions must hold everywhere.}
        \label{fig:general-a-allspace}
    \end{subfigure}
    \hfill
    \begin{subfigure}[t]{0.31\linewidth}
        \centering
        \vspace{0pt}
        \includegraphics[width=\linewidth]{images/figure_2b.pdf}
        \caption{\citet{rubinstein2026accuracy}: it suffices to check conditions
        in a ``sufficiently large'' ellipsoid.}
        \label{fig:general-a-ellipsoid}
    \end{subfigure}
    \hfill
    \begin{subfigure}[t]{0.31\linewidth}
        \centering
        \vspace{0pt}
        \includegraphics[width=\linewidth]{images/figure_2c.pdf}
        \caption{Theorem~\ref{thm:general_intro}: any sufficiently large convex
        set suffices.}
        \label{fig:general-a-balanced}
    \end{subfigure}
    \caption[Progress in theoretical analyses of approximate unlearning.]{%
        Progress in theoretical analyses of approximate unlearning algorithms,
        inspired by Figure~2 of \citet{rubinstein2026accuracy}.
        \textbf{(a)} Until recently, the leading analyses
        \citep{TODO-global-analyses} required global convexity and smoothness
        conditions on the loss, necessitating a very large regularization
        coefficient for their guarantees to apply.
        \textbf{(b)} \citet{rubinstein2026accuracy} showed that it suffices for
        these conditions to hold within a relatively small ellipsoid.
        This improved on the previous state of the art because these local
        conditions hold for ``well-behaved'' regression problems. However,
        turning this result into an unlearning algorithm requires a
        corresponding certification algorithm that verifies the local
        conditions for a specific dataset.
        \textbf{(c)} In Theorem~\ref{thm:general_intro}, we show that it suffices
        to verify the same conditions in any ``sufficiently large'' convex
        set $\mcA$. This seemingly minor difference is significant for our
        \unerring{} algorithm, which applies Theorem~\ref{thm:general_intro}
        to the intersection of an ellipsoid and a parallelepiped, and
        requires a fundamentally different proof technique.
        All three panels depict the gradient-descent field, $-\nabla L(\theta)$.
    }
    \label{fig:general-a-comparison}
\end{figure}
